(1) 封闭量子系统初始归一态|psi0>，时间无关自伴哈密顿量H且初态具有有限<H²>。展开不测量的短时生存振幅A(t)=<psi0|exp(-iHt/hbar)|psi0>到t²，求生存概率。
(2) 每隔tau=t/N对初态投影测量并条件化为“仍在初态”，求N次全部成功概率P_N(t)，证明N趋无穷时的Zeno极限。
(3) 在有限N、短tau范围求有效指数衰减率和Zeno时间，比较连续未观测演化的短时二次衰减与常数Markov衰减率模型。
(4) 解释为什么真实测量需具有带宽、耦合能量和环境效应，以及在能量方差发散或非理想POVM情况下二次律和冻结极限可能失效。
